\chapter{Getting started} \label{ch1}

Setting up a new thesis only requires editing the preamble of
\texttt{main.tex}: choose the degree programme, then fill in the thesis
data. This chapter describes both steps, and the commands that organise the
document.

\section{Choosing the degree programme} \label{sec:programme}

The degree programme is the main option of the style:
\begin{Verbatim}
\usepackage[programme=dse]{slimthesis}
\end{Verbatim}
It determines the cover (logo, faculty and course), the language of the
document and all the fixed labels, as summarised in
Table~\ref{tab:programmes}.

\begin{table}[htbp]
  \centering
  \caption{Values of the \texttt{programme} option.}
  \label{tab:programmes}
  \begin{tabular}{llll}
    \toprule
    Value & Degree & Programme & Language \\
    \midrule
    \texttt{inf} (default) & Bachelor's & Informatica & Italian \\
    \texttt{inf-lm} & Master's & Informatica & Italian \\
    \texttt{cs}     & Master's & Computer Science & English \\
    \texttt{dse}    & Master's & Data Science for Economics & English \\
    \texttt{dseh}   & Master's & Data Science for Economics and Health & English \\
    \bottomrule
  \end{tabular}
\end{table}

The language affects chapter labels (``Capitolo''/``Chapter''), the titles of
the table of contents and of the bibliography, the names of theorem-like
environments (Section~\ref{sec:theorems}) and the labels on the cover. If
some passages are written in another language, add it to babel
\emph{before} loading the style, and switch to it where needed:
\begin{Verbatim}
\PassOptionsToPackage{italian}{babel}
\usepackage[programme=dse]{slimthesis}
...
\foreignlanguage{italian}{Questa frase è in italiano.}
\end{Verbatim}

The other options of the style are rarely needed:
\begin{description}
  \item[\texttt{madeinslim=false}] omits the closing ``Made in SLIM'' page;
  \item[\texttt{pdfa=false}] produces an ordinary PDF file instead of a
    PDF/A-2b one; links are then blue on screen and black in print
    (Section~\ref{ch3});
  \item[\texttt{fontdir=\textrm{\emph{dir}}}] directory containing the
    Figtree fonts (default \nolinkurl{fonts/Figtree/});
  \item[\texttt{mathfont=\textrm{\emph{name}}}] the OpenType math font
    (default \texttt{STIX Two Math}). Any installed math font can be used,
    for instance \texttt{Libertinus Math}, \texttt{Latin Modern Math} or
    \texttt{TeX Gyre Pagella Math}; it is automatically scaled to match the
    size of the text font;
  \item[\texttt{logodir=\textrm{\emph{dir}}}] directory containing the cover
    logos (default \nolinkurl{images/coverPage/});
  \item[\texttt{licensedir=\textrm{\emph{dir}}}] directory containing the
    Creative Commons buttons (default \nolinkurl{images/license/}).
\end{description}
Options are separated by commas; a font name containing spaces needs no
quotes:
\begin{Verbatim}
\usepackage[programme=cs, mathfont=Libertinus Math]{slimthesis}
\end{Verbatim}

\section{Thesis data and cover} \label{sec:data}

The cover is produced by \cs{maketitle}, using the data given in the
preamble:
\begin{Verbatim}
\title{Learning Curves of Ensemble Classifiers}
\author{Ada Lovelace}
\academicyear{2025/26}
\studentid{123456}
\keywords{ensemble methods, random forests, learning curves}
\supervisor{Prof.~Super Visor}
\cosupervisor{Prof.~Cosu Pervisor}
\cosupervisor{Dr.~Other Supervisor}
\end{Verbatim}
\cs{academicyear} only takes the years: ``Anno Accademico'' or ``Academic
Year'' is added according to the language. \cs{studentid} is optional; when
present, the student ID is printed below the author's name. \cs{keywords}
is optional too: the keywords are not printed, but stored in the metadata of
the PDF file, together with title, author and degree course, so that search
engines and repositories can index the thesis.
\cs{cosupervisor} can be repeated, or omitted: the label is singular or
plural depending on the number of co-supervisors, and disappears if there
are none. Write \verb|Prof.~| and \verb|Dr.~| with a tie, so that no
sentence-ending space is added after the abbreviation and title and name
stay on the same line.

For the Informatica programmes the cover shows the official logo, which
already contains faculty and course. For the other programmes the faculty
and course are typeset under the university seal; interfaculty programmes
(DSE and DSEH) list both faculties, at the same size. These lines can be
changed with
\begin{Verbatim}
\faculty{Facoltà di Scienze e Tecnologie}
\course{Master's Degree in Data Science\\for Economics and Health}
\end{Verbatim}
The first \cs{faculty} replaces the default faculties, further calls add
lines. A course name that is too long is wrapped automatically; use
\verb|\\| to choose where it breaks.

\subsection{License} \label{sec:license}

A thesis can be released under a license, which tells readers what they are
allowed to do with it. The license is printed at the bottom of the page
following the cover. For Creative Commons licenses, give their type:
\begin{Verbatim}
\license{by-sa}
\end{Verbatim}
The available types are listed in Table~\ref{tab:licenses}; an optional
argument selects a version other than 4.0, as in \verb|\license[3.0]{by}|.
The page shows the official button of the license, a copyright line with
the current year and the author's name (except for CC0, which waives
copyright), and the name and address of the license, in the language of
the document. Any other license or notice can be given as free text:
\begin{Verbatim}
\customlicense{\textcopyright\ 2026 Ada Lovelace. All rights reserved.}
\end{Verbatim}
Without \cs{license} or \cs{customlicense}, the page is omitted.

\begin{table}[htbp]
  \centering
  \caption{Creative Commons licenses available through \cs{license}.}
  \label{tab:licenses}
  \begin{tabular}{ll}
    \toprule
    Type & License \\
    \midrule
    \texttt{by}       & Attribution \\
    \texttt{by-sa}    & Attribution-ShareAlike \\
    \texttt{by-nd}    & Attribution-NoDerivatives \\
    \texttt{by-nc}    & Attribution-NonCommercial \\
    \texttt{by-nc-sa} & Attribution-NonCommercial-ShareAlike \\
    \texttt{by-nc-nd} & Attribution-NonCommercial-NoDerivatives \\
    \texttt{zero}     & CC0 (public domain dedication) \\
    \bottomrule
  \end{tabular}
\end{table}

\section{Structure of the document} \label{sec:structure}

The body of \texttt{main.tex} follows the usual structure of a book:
\begin{Verbatim}
\begin{document}
\frontmatter
\maketitle
\dedication{To someone special\ldots}

\epigraphpage{A quotation \\
one of those that make you think.}{--- A. Uthor}

\chapter*{Acknowledgements}

Thank the people who helped you during your studies and while writing
the thesis.

Or simply strip out this page.

\tableofcontents
\listoffigures
\listoftables
\listofalgorithms
\begin{listofsymbols}
  \item[$\R$]                 the set of real numbers
  \item[$d$]                  dimension of the data space
  \item[$n$]                  number of points in a dataset
  \item[$\bx, \by$]           vectors of $\R^d$, written in bold italic
  \item[$\bx_i$]              the $i$-th point of a dataset
  \item[$\lVert \bx \rVert$]  Euclidean norm of $\bx$
  \item[$\symbfit{\mu}_j$]    centroid of the $j$-th cluster
  \item[$\sigma(z)$]          logistic function, $1 / (1 + e^{-z})$
  \item[$\ell(y, z)$]         loss of prediction $z$ when the label is $y$
  \item[$\F_1$]               harmonic mean of precision and recall
\end{listofsymbols}


\mainmatter
\chapter*{Introduction}

This document is at the same time the documentation and an example of
\texttt{slimthesis}, a \LaTeX{} style for bachelor's and master's theses
supervised at the SLIM laboratory of the Università degli Studi di Milano.
Its source is the starting point for a new thesis: replace the text of the
chapters with your own, keep the structure, and the style takes care of the
cover page, the typography and the bibliography.

All the layout is in a single file, \texttt{slimthesis.sty}. The document
itself (\texttt{main.tex} and the files it includes) only contains the thesis
data and the text, so that the style can be updated without touching your
work, and your work never needs to touch the style.

\paragraph{Requirements}

The style relies on OpenType fonts, therefore it must be compiled with
\emph{LuaLaTeX} or \emph{XeLaTeX}; pdf\LaTeX{} is not supported. Besides a
recent \TeX{} Live distribution, you need:
\begin{itemize}
  \item the \emph{STIX Two Math} font, used for mathematics (the Fedora
    package is \nolinkurl{stix-fonts});
  \item the \emph{Libertinus} fonts, used for text (the
    Fedora package is \nolinkurl{texlive-libertinus-fonts});
  \item the \emph{Figtree} font, used for headings, which is shipped with the
    template in \nolinkurl{fonts/Figtree/};
  \item a few \LaTeX{} packages that minimal installations may lack, such as
    \texttt{apptools}, \texttt{ocgx2} and \texttt{pdfx} (on Fedora:
    \nolinkurl{texlive-apptools}, \nolinkurl{texlive-ocgx2} and
    \nolinkurl{texlive-pdfx});
  \item \emph{biber}, which processes the bibliography.
\end{itemize}

\paragraph{Compiling}

A complete build runs LuaLaTeX, then biber, then LuaLaTeX twice more, so
that the bibliography, the table of contents and all cross-references are
resolved:
\begin{Verbatim}
lualatex main && biber main && lualatex main && lualatex main
\end{Verbatim}
Alternatively, \texttt{latexmk} runs all the required steps automatically;
the file \texttt{latexmkrc} shipped with the template tells it to use
LuaLaTeX and biber, so that it is enough to type
\begin{Verbatim}
latexmk main
\end{Verbatim}
and \verb|latexmk -c| removes the auxiliary files. XeLaTeX can be used in
the same way, replacing \texttt{lualatex} with \texttt{xelatex}.

\paragraph{PDF/A}

The thesis is produced as a PDF/A-2b file, the archival format usually
required when the final version is deposited, through the \texttt{pdfx}
package. The title, the author, the degree course and the keywords given
with \cs{keywords} are stored in the metadata of the file; the style writes
them in the file \texttt{main.xmpdata}, read by \texttt{pdfx}, which is
regenerated at each compilation and should not be edited. Conformance can be
checked with the free validator veraPDF (\url{https://verapdf.org}). The
option \texttt{pdfa=false} produces an ordinary PDF file instead
(Section~\ref{sec:programme}).

\paragraph{Files}

\begin{description}
  \item[\texttt{main.tex}] the main file: programme, thesis data and the list
    of chapters;
  \item[\texttt{slimthesis.sty}] the style; normally it does not need to be
    modified;
  \item[\texttt{frontmatter/}] dedication, epigraph and acknowledgements;
  \item[\texttt{chapters/}] one file per chapter;
  \item[\texttt{bibliography.bib}] the bibliographic references;
  \item[\texttt{style/}] your own macros (\texttt{symbols.tex}) and the
    colours used by the figures (\texttt{colors.tex});
  \item[\texttt{images/}] figures, the logos used on the cover
    (\texttt{images/coverPage/}) and the Creative Commons buttons
    (\texttt{images/license/});
  \item[\texttt{fonts/}] the Figtree font.
\end{description}

Chapter~\ref{ch1} explains how to set up the thesis, Chapter~\ref{ch2} how
to write it, Chapter~\ref{ch3} how to cite and link, and the last chapter
lists a few checks to do before submitting.

\chapter{Getting started} \label{ch1}

Setting up a new thesis only requires editing the preamble of
\texttt{main.tex}: choose the degree programme, then fill in the thesis
data. This chapter describes both steps, and the commands that organise the
document.

\section{Choosing the degree programme} \label{sec:programme}

The degree programme is the main option of the style:
\begin{Verbatim}
\usepackage[programme=dse]{slimthesis}
\end{Verbatim}
It determines the cover (logo, faculty and course), the language of the
document and all the fixed labels, as summarised in
Table~\ref{tab:programmes}.

\begin{table}[htbp]
  \centering
  \caption{Values of the \texttt{programme} option.}
  \label{tab:programmes}
  \begin{tabular}{llll}
    \toprule
    Value & Degree & Programme & Language \\
    \midrule
    \texttt{inf} (default) & Bachelor's & Informatica & Italian \\
    \texttt{inf-lm} & Master's & Informatica & Italian \\
    \texttt{cs}     & Master's & Computer Science & English \\
    \texttt{dse}    & Master's & Data Science for Economics & English \\
    \texttt{dseh}   & Master's & Data Science for Economics and Health & English \\
    \bottomrule
  \end{tabular}
\end{table}

The language affects chapter labels (``Capitolo''/``Chapter''), the titles of
the table of contents and of the bibliography, the names of theorem-like
environments (Section~\ref{sec:theorems}) and the labels on the cover. If
some passages are written in another language, add it to babel
\emph{before} loading the style, and switch to it where needed:
\begin{Verbatim}
\PassOptionsToPackage{italian}{babel}
\usepackage[programme=dse]{slimthesis}
...
\foreignlanguage{italian}{Questa frase è in italiano.}
\end{Verbatim}

The other options of the style are rarely needed:
\begin{description}
  \item[\texttt{madeinslim=false}] omits the closing ``Made in SLIM'' page;
  \item[\texttt{pdfa=false}] produces an ordinary PDF file instead of a
    PDF/A-2b one; links are then blue on screen and black in print
    (Section~\ref{ch3});
  \item[\texttt{fontdir=\textrm{\emph{dir}}}] directory containing the
    Figtree fonts (default \nolinkurl{fonts/Figtree/});
  \item[\texttt{mathfont=\textrm{\emph{name}}}] the OpenType math font
    (default \texttt{STIX Two Math}). Any installed math font can be used,
    for instance \texttt{Libertinus Math}, \texttt{Latin Modern Math} or
    \texttt{TeX Gyre Pagella Math}; it is automatically scaled to match the
    size of the text font;
  \item[\texttt{logodir=\textrm{\emph{dir}}}] directory containing the cover
    logos (default \nolinkurl{images/coverPage/});
  \item[\texttt{licensedir=\textrm{\emph{dir}}}] directory containing the
    Creative Commons buttons (default \nolinkurl{images/license/}).
\end{description}
Options are separated by commas; a font name containing spaces needs no
quotes:
\begin{Verbatim}
\usepackage[programme=cs, mathfont=Libertinus Math]{slimthesis}
\end{Verbatim}

\section{Thesis data and cover} \label{sec:data}

The cover is produced by \cs{maketitle}, using the data given in the
preamble:
\begin{Verbatim}
\title{Learning Curves of Ensemble Classifiers}
\author{Ada Lovelace}
\academicyear{2025/26}
\studentid{123456}
\keywords{ensemble methods, random forests, learning curves}
\supervisor{Prof.~Super Visor}
\cosupervisor{Prof.~Cosu Pervisor}
\cosupervisor{Dr.~Other Supervisor}
\end{Verbatim}
\cs{academicyear} only takes the years: ``Anno Accademico'' or ``Academic
Year'' is added according to the language. \cs{studentid} is optional; when
present, the student ID is printed below the author's name. \cs{keywords}
is optional too: the keywords are not printed, but stored in the metadata of
the PDF file, together with title, author and degree course, so that search
engines and repositories can index the thesis.
\cs{cosupervisor} can be repeated, or omitted: the label is singular or
plural depending on the number of co-supervisors, and disappears if there
are none. Write \verb|Prof.~| and \verb|Dr.~| with a tie, so that no
sentence-ending space is added after the abbreviation and title and name
stay on the same line.

For the Informatica programmes the cover shows the official logo, which
already contains faculty and course. For the other programmes the faculty
and course are typeset under the university seal; interfaculty programmes
(DSE and DSEH) list both faculties, at the same size. These lines can be
changed with
\begin{Verbatim}
\faculty{Facoltà di Scienze e Tecnologie}
\course{Master's Degree in Data Science\\for Economics and Health}
\end{Verbatim}
The first \cs{faculty} replaces the default faculties, further calls add
lines. A course name that is too long is wrapped automatically; use
\verb|\\| to choose where it breaks.

\subsection{License} \label{sec:license}

A thesis can be released under a license, which tells readers what they are
allowed to do with it. The license is printed at the bottom of the page
following the cover. For Creative Commons licenses, give their type:
\begin{Verbatim}
\license{by-sa}
\end{Verbatim}
The available types are listed in Table~\ref{tab:licenses}; an optional
argument selects a version other than 4.0, as in \verb|\license[3.0]{by}|.
The page shows the official button of the license, a copyright line with
the current year and the author's name (except for CC0, which waives
copyright), and the name and address of the license, in the language of
the document. Any other license or notice can be given as free text:
\begin{Verbatim}
\customlicense{\textcopyright\ 2026 Ada Lovelace. All rights reserved.}
\end{Verbatim}
Without \cs{license} or \cs{customlicense}, the page is omitted.

\begin{table}[htbp]
  \centering
  \caption{Creative Commons licenses available through \cs{license}.}
  \label{tab:licenses}
  \begin{tabular}{ll}
    \toprule
    Type & License \\
    \midrule
    \texttt{by}       & Attribution \\
    \texttt{by-sa}    & Attribution-ShareAlike \\
    \texttt{by-nd}    & Attribution-NoDerivatives \\
    \texttt{by-nc}    & Attribution-NonCommercial \\
    \texttt{by-nc-sa} & Attribution-NonCommercial-ShareAlike \\
    \texttt{by-nc-nd} & Attribution-NonCommercial-NoDerivatives \\
    \texttt{zero}     & CC0 (public domain dedication) \\
    \bottomrule
  \end{tabular}
\end{table}

\section{Structure of the document} \label{sec:structure}

The body of \texttt{main.tex} follows the usual structure of a book:
\begin{Verbatim}
\begin{document}
\frontmatter
\maketitle
\dedication{To someone special\ldots}

\epigraphpage{A quotation \\
one of those that make you think.}{--- A. Uthor}

\chapter*{Acknowledgements}

Thank the people who helped you during your studies and while writing
the thesis.

Or simply strip out this page.

\tableofcontents
\listoffigures
\listoftables
\listofalgorithms
\begin{listofsymbols}
  \item[$\R$]                 the set of real numbers
  \item[$d$]                  dimension of the data space
  \item[$n$]                  number of points in a dataset
  \item[$\bx, \by$]           vectors of $\R^d$, written in bold italic
  \item[$\bx_i$]              the $i$-th point of a dataset
  \item[$\lVert \bx \rVert$]  Euclidean norm of $\bx$
  \item[$\symbfit{\mu}_j$]    centroid of the $j$-th cluster
  \item[$\sigma(z)$]          logistic function, $1 / (1 + e^{-z})$
  \item[$\ell(y, z)$]         loss of prediction $z$ when the label is $y$
  \item[$\F_1$]               harmonic mean of precision and recall
\end{listofsymbols}


\mainmatter
\chapter*{Introduction}

This document is at the same time the documentation and an example of
\texttt{slimthesis}, a \LaTeX{} style for bachelor's and master's theses
supervised at the SLIM laboratory of the Università degli Studi di Milano.
Its source is the starting point for a new thesis: replace the text of the
chapters with your own, keep the structure, and the style takes care of the
cover page, the typography and the bibliography.

All the layout is in a single file, \texttt{slimthesis.sty}. The document
itself (\texttt{main.tex} and the files it includes) only contains the thesis
data and the text, so that the style can be updated without touching your
work, and your work never needs to touch the style.

\paragraph{Requirements}

The style relies on OpenType fonts, therefore it must be compiled with
\emph{LuaLaTeX} or \emph{XeLaTeX}; pdf\LaTeX{} is not supported. Besides a
recent \TeX{} Live distribution, you need:
\begin{itemize}
  \item the \emph{STIX Two Math} font, used for mathematics (the Fedora
    package is \nolinkurl{stix-fonts});
  \item the \emph{Libertinus} fonts, used for text (the
    Fedora package is \nolinkurl{texlive-libertinus-fonts});
  \item the \emph{Figtree} font, used for headings, which is shipped with the
    template in \nolinkurl{fonts/Figtree/};
  \item a few \LaTeX{} packages that minimal installations may lack, such as
    \texttt{apptools}, \texttt{ocgx2} and \texttt{pdfx} (on Fedora:
    \nolinkurl{texlive-apptools}, \nolinkurl{texlive-ocgx2} and
    \nolinkurl{texlive-pdfx});
  \item \emph{biber}, which processes the bibliography.
\end{itemize}

\paragraph{Compiling}

A complete build runs LuaLaTeX, then biber, then LuaLaTeX twice more, so
that the bibliography, the table of contents and all cross-references are
resolved:
\begin{Verbatim}
lualatex main && biber main && lualatex main && lualatex main
\end{Verbatim}
Alternatively, \texttt{latexmk} runs all the required steps automatically;
the file \texttt{latexmkrc} shipped with the template tells it to use
LuaLaTeX and biber, so that it is enough to type
\begin{Verbatim}
latexmk main
\end{Verbatim}
and \verb|latexmk -c| removes the auxiliary files. XeLaTeX can be used in
the same way, replacing \texttt{lualatex} with \texttt{xelatex}.

\paragraph{PDF/A}

The thesis is produced as a PDF/A-2b file, the archival format usually
required when the final version is deposited, through the \texttt{pdfx}
package. The title, the author, the degree course and the keywords given
with \cs{keywords} are stored in the metadata of the file; the style writes
them in the file \texttt{main.xmpdata}, read by \texttt{pdfx}, which is
regenerated at each compilation and should not be edited. Conformance can be
checked with the free validator veraPDF (\url{https://verapdf.org}). The
option \texttt{pdfa=false} produces an ordinary PDF file instead
(Section~\ref{sec:programme}).

\paragraph{Files}

\begin{description}
  \item[\texttt{main.tex}] the main file: programme, thesis data and the list
    of chapters;
  \item[\texttt{slimthesis.sty}] the style; normally it does not need to be
    modified;
  \item[\texttt{frontmatter/}] dedication, epigraph and acknowledgements;
  \item[\texttt{chapters/}] one file per chapter;
  \item[\texttt{bibliography.bib}] the bibliographic references;
  \item[\texttt{style/}] your own macros (\texttt{symbols.tex}) and the
    colours used by the figures (\texttt{colors.tex});
  \item[\texttt{images/}] figures, the logos used on the cover
    (\texttt{images/coverPage/}) and the Creative Commons buttons
    (\texttt{images/license/});
  \item[\texttt{fonts/}] the Figtree font.
\end{description}

Chapter~\ref{ch1} explains how to set up the thesis, Chapter~\ref{ch2} how
to write it, Chapter~\ref{ch3} how to cite and link, and the last chapter
lists a few checks to do before submitting.

\chapter{Getting started} \label{ch1}

Setting up a new thesis only requires editing the preamble of
\texttt{main.tex}: choose the degree programme, then fill in the thesis
data. This chapter describes both steps, and the commands that organise the
document.

\section{Choosing the degree programme} \label{sec:programme}

The degree programme is the main option of the style:
\begin{Verbatim}
\usepackage[programme=dse]{slimthesis}
\end{Verbatim}
It determines the cover (logo, faculty and course), the language of the
document and all the fixed labels, as summarised in
Table~\ref{tab:programmes}.

\begin{table}[htbp]
  \centering
  \caption{Values of the \texttt{programme} option.}
  \label{tab:programmes}
  \begin{tabular}{llll}
    \toprule
    Value & Degree & Programme & Language \\
    \midrule
    \texttt{inf} (default) & Bachelor's & Informatica & Italian \\
    \texttt{inf-lm} & Master's & Informatica & Italian \\
    \texttt{cs}     & Master's & Computer Science & English \\
    \texttt{dse}    & Master's & Data Science for Economics & English \\
    \texttt{dseh}   & Master's & Data Science for Economics and Health & English \\
    \bottomrule
  \end{tabular}
\end{table}

The language affects chapter labels (``Capitolo''/``Chapter''), the titles of
the table of contents and of the bibliography, the names of theorem-like
environments (Section~\ref{sec:theorems}) and the labels on the cover. If
some passages are written in another language, add it to babel
\emph{before} loading the style, and switch to it where needed:
\begin{Verbatim}
\PassOptionsToPackage{italian}{babel}
\usepackage[programme=dse]{slimthesis}
...
\foreignlanguage{italian}{Questa frase è in italiano.}
\end{Verbatim}

The other options of the style are rarely needed:
\begin{description}
  \item[\texttt{madeinslim=false}] omits the closing ``Made in SLIM'' page;
  \item[\texttt{pdfa=false}] produces an ordinary PDF file instead of a
    PDF/A-2b one; links are then blue on screen and black in print
    (Section~\ref{ch3});
  \item[\texttt{fontdir=\textrm{\emph{dir}}}] directory containing the
    Figtree fonts (default \nolinkurl{fonts/Figtree/});
  \item[\texttt{mathfont=\textrm{\emph{name}}}] the OpenType math font
    (default \texttt{STIX Two Math}). Any installed math font can be used,
    for instance \texttt{Libertinus Math}, \texttt{Latin Modern Math} or
    \texttt{TeX Gyre Pagella Math}; it is automatically scaled to match the
    size of the text font;
  \item[\texttt{logodir=\textrm{\emph{dir}}}] directory containing the cover
    logos (default \nolinkurl{images/coverPage/});
  \item[\texttt{licensedir=\textrm{\emph{dir}}}] directory containing the
    Creative Commons buttons (default \nolinkurl{images/license/}).
\end{description}
Options are separated by commas; a font name containing spaces needs no
quotes:
\begin{Verbatim}
\usepackage[programme=cs, mathfont=Libertinus Math]{slimthesis}
\end{Verbatim}

\section{Thesis data and cover} \label{sec:data}

The cover is produced by \cs{maketitle}, using the data given in the
preamble:
\begin{Verbatim}
\title{Learning Curves of Ensemble Classifiers}
\author{Ada Lovelace}
\academicyear{2025/26}
\studentid{123456}
\keywords{ensemble methods, random forests, learning curves}
\supervisor{Prof.~Super Visor}
\cosupervisor{Prof.~Cosu Pervisor}
\cosupervisor{Dr.~Other Supervisor}
\end{Verbatim}
\cs{academicyear} only takes the years: ``Anno Accademico'' or ``Academic
Year'' is added according to the language. \cs{studentid} is optional; when
present, the student ID is printed below the author's name. \cs{keywords}
is optional too: the keywords are not printed, but stored in the metadata of
the PDF file, together with title, author and degree course, so that search
engines and repositories can index the thesis.
\cs{cosupervisor} can be repeated, or omitted: the label is singular or
plural depending on the number of co-supervisors, and disappears if there
are none. Write \verb|Prof.~| and \verb|Dr.~| with a tie, so that no
sentence-ending space is added after the abbreviation and title and name
stay on the same line.

For the Informatica programmes the cover shows the official logo, which
already contains faculty and course. For the other programmes the faculty
and course are typeset under the university seal; interfaculty programmes
(DSE and DSEH) list both faculties, at the same size. These lines can be
changed with
\begin{Verbatim}
\faculty{Facoltà di Scienze e Tecnologie}
\course{Master's Degree in Data Science\\for Economics and Health}
\end{Verbatim}
The first \cs{faculty} replaces the default faculties, further calls add
lines. A course name that is too long is wrapped automatically; use
\verb|\\| to choose where it breaks.

\subsection{License} \label{sec:license}

A thesis can be released under a license, which tells readers what they are
allowed to do with it. The license is printed at the bottom of the page
following the cover. For Creative Commons licenses, give their type:
\begin{Verbatim}
\license{by-sa}
\end{Verbatim}
The available types are listed in Table~\ref{tab:licenses}; an optional
argument selects a version other than 4.0, as in \verb|\license[3.0]{by}|.
The page shows the official button of the license, a copyright line with
the current year and the author's name (except for CC0, which waives
copyright), and the name and address of the license, in the language of
the document. Any other license or notice can be given as free text:
\begin{Verbatim}
\customlicense{\textcopyright\ 2026 Ada Lovelace. All rights reserved.}
\end{Verbatim}
Without \cs{license} or \cs{customlicense}, the page is omitted.

\begin{table}[htbp]
  \centering
  \caption{Creative Commons licenses available through \cs{license}.}
  \label{tab:licenses}
  \begin{tabular}{ll}
    \toprule
    Type & License \\
    \midrule
    \texttt{by}       & Attribution \\
    \texttt{by-sa}    & Attribution-ShareAlike \\
    \texttt{by-nd}    & Attribution-NoDerivatives \\
    \texttt{by-nc}    & Attribution-NonCommercial \\
    \texttt{by-nc-sa} & Attribution-NonCommercial-ShareAlike \\
    \texttt{by-nc-nd} & Attribution-NonCommercial-NoDerivatives \\
    \texttt{zero}     & CC0 (public domain dedication) \\
    \bottomrule
  \end{tabular}
\end{table}

\section{Structure of the document} \label{sec:structure}

The body of \texttt{main.tex} follows the usual structure of a book:
\begin{Verbatim}
\begin{document}
\frontmatter
\maketitle
\dedication{To someone special\ldots}

\epigraphpage{A quotation \\
one of those that make you think.}{--- A. Uthor}

\chapter*{Acknowledgements}

Thank the people who helped you during your studies and while writing
the thesis.

Or simply strip out this page.

\tableofcontents
\listoffigures
\listoftables
\listofalgorithms
\begin{listofsymbols}
  \item[$\R$]                 the set of real numbers
  \item[$d$]                  dimension of the data space
  \item[$n$]                  number of points in a dataset
  \item[$\bx, \by$]           vectors of $\R^d$, written in bold italic
  \item[$\bx_i$]              the $i$-th point of a dataset
  \item[$\lVert \bx \rVert$]  Euclidean norm of $\bx$
  \item[$\symbfit{\mu}_j$]    centroid of the $j$-th cluster
  \item[$\sigma(z)$]          logistic function, $1 / (1 + e^{-z})$
  \item[$\ell(y, z)$]         loss of prediction $z$ when the label is $y$
  \item[$\F_1$]               harmonic mean of precision and recall
\end{listofsymbols}


\mainmatter
\chapter*{Introduction}

This document is at the same time the documentation and an example of
\texttt{slimthesis}, a \LaTeX{} style for bachelor's and master's theses
supervised at the SLIM laboratory of the Università degli Studi di Milano.
Its source is the starting point for a new thesis: replace the text of the
chapters with your own, keep the structure, and the style takes care of the
cover page, the typography and the bibliography.

All the layout is in a single file, \texttt{slimthesis.sty}. The document
itself (\texttt{main.tex} and the files it includes) only contains the thesis
data and the text, so that the style can be updated without touching your
work, and your work never needs to touch the style.

\paragraph{Requirements}

The style relies on OpenType fonts, therefore it must be compiled with
\emph{LuaLaTeX} or \emph{XeLaTeX}; pdf\LaTeX{} is not supported. Besides a
recent \TeX{} Live distribution, you need:
\begin{itemize}
  \item the \emph{STIX Two Math} font, used for mathematics (the Fedora
    package is \nolinkurl{stix-fonts});
  \item the \emph{Libertinus} fonts, used for text (the
    Fedora package is \nolinkurl{texlive-libertinus-fonts});
  \item the \emph{Figtree} font, used for headings, which is shipped with the
    template in \nolinkurl{fonts/Figtree/};
  \item a few \LaTeX{} packages that minimal installations may lack, such as
    \texttt{apptools}, \texttt{ocgx2} and \texttt{pdfx} (on Fedora:
    \nolinkurl{texlive-apptools}, \nolinkurl{texlive-ocgx2} and
    \nolinkurl{texlive-pdfx});
  \item \emph{biber}, which processes the bibliography.
\end{itemize}

\paragraph{Compiling}

A complete build runs LuaLaTeX, then biber, then LuaLaTeX twice more, so
that the bibliography, the table of contents and all cross-references are
resolved:
\begin{Verbatim}
lualatex main && biber main && lualatex main && lualatex main
\end{Verbatim}
Alternatively, \texttt{latexmk} runs all the required steps automatically;
the file \texttt{latexmkrc} shipped with the template tells it to use
LuaLaTeX and biber, so that it is enough to type
\begin{Verbatim}
latexmk main
\end{Verbatim}
and \verb|latexmk -c| removes the auxiliary files. XeLaTeX can be used in
the same way, replacing \texttt{lualatex} with \texttt{xelatex}.

\paragraph{PDF/A}

The thesis is produced as a PDF/A-2b file, the archival format usually
required when the final version is deposited, through the \texttt{pdfx}
package. The title, the author, the degree course and the keywords given
with \cs{keywords} are stored in the metadata of the file; the style writes
them in the file \texttt{main.xmpdata}, read by \texttt{pdfx}, which is
regenerated at each compilation and should not be edited. Conformance can be
checked with the free validator veraPDF (\url{https://verapdf.org}). The
option \texttt{pdfa=false} produces an ordinary PDF file instead
(Section~\ref{sec:programme}).

\paragraph{Files}

\begin{description}
  \item[\texttt{main.tex}] the main file: programme, thesis data and the list
    of chapters;
  \item[\texttt{slimthesis.sty}] the style; normally it does not need to be
    modified;
  \item[\texttt{frontmatter/}] dedication, epigraph and acknowledgements;
  \item[\texttt{chapters/}] one file per chapter;
  \item[\texttt{bibliography.bib}] the bibliographic references;
  \item[\texttt{style/}] your own macros (\texttt{symbols.tex}) and the
    colours used by the figures (\texttt{colors.tex});
  \item[\texttt{images/}] figures, the logos used on the cover
    (\texttt{images/coverPage/}) and the Creative Commons buttons
    (\texttt{images/license/});
  \item[\texttt{fonts/}] the Figtree font.
\end{description}

Chapter~\ref{ch1} explains how to set up the thesis, Chapter~\ref{ch2} how
to write it, Chapter~\ref{ch3} how to cite and link, and the last chapter
lists a few checks to do before submitting.

\chapter{Getting started} \label{ch1}

Setting up a new thesis only requires editing the preamble of
\texttt{main.tex}: choose the degree programme, then fill in the thesis
data. This chapter describes both steps, and the commands that organise the
document.

\section{Choosing the degree programme} \label{sec:programme}

The degree programme is the main option of the style:
\begin{Verbatim}
\usepackage[programme=dse]{slimthesis}
\end{Verbatim}
It determines the cover (logo, faculty and course), the language of the
document and all the fixed labels, as summarised in
Table~\ref{tab:programmes}.

\begin{table}[htbp]
  \centering
  \caption{Values of the \texttt{programme} option.}
  \label{tab:programmes}
  \begin{tabular}{llll}
    \toprule
    Value & Degree & Programme & Language \\
    \midrule
    \texttt{inf} (default) & Bachelor's & Informatica & Italian \\
    \texttt{inf-lm} & Master's & Informatica & Italian \\
    \texttt{cs}     & Master's & Computer Science & English \\
    \texttt{dse}    & Master's & Data Science for Economics & English \\
    \texttt{dseh}   & Master's & Data Science for Economics and Health & English \\
    \bottomrule
  \end{tabular}
\end{table}

The language affects chapter labels (``Capitolo''/``Chapter''), the titles of
the table of contents and of the bibliography, the names of theorem-like
environments (Section~\ref{sec:theorems}) and the labels on the cover. If
some passages are written in another language, add it to babel
\emph{before} loading the style, and switch to it where needed:
\begin{Verbatim}
\PassOptionsToPackage{italian}{babel}
\usepackage[programme=dse]{slimthesis}
...
\foreignlanguage{italian}{Questa frase è in italiano.}
\end{Verbatim}

The other options of the style are rarely needed:
\begin{description}
  \item[\texttt{madeinslim=false}] omits the closing ``Made in SLIM'' page;
  \item[\texttt{pdfa=false}] produces an ordinary PDF file instead of a
    PDF/A-2b one; links are then blue on screen and black in print
    (Section~\ref{ch3});
  \item[\texttt{fontdir=\textrm{\emph{dir}}}] directory containing the
    Figtree fonts (default \nolinkurl{fonts/Figtree/});
  \item[\texttt{mathfont=\textrm{\emph{name}}}] the OpenType math font
    (default \texttt{STIX Two Math}). Any installed math font can be used,
    for instance \texttt{Libertinus Math}, \texttt{Latin Modern Math} or
    \texttt{TeX Gyre Pagella Math}; it is automatically scaled to match the
    size of the text font;
  \item[\texttt{logodir=\textrm{\emph{dir}}}] directory containing the cover
    logos (default \nolinkurl{images/coverPage/});
  \item[\texttt{licensedir=\textrm{\emph{dir}}}] directory containing the
    Creative Commons buttons (default \nolinkurl{images/license/}).
\end{description}
Options are separated by commas; a font name containing spaces needs no
quotes:
\begin{Verbatim}
\usepackage[programme=cs, mathfont=Libertinus Math]{slimthesis}
\end{Verbatim}

\section{Thesis data and cover} \label{sec:data}

The cover is produced by \cs{maketitle}, using the data given in the
preamble:
\begin{Verbatim}
\title{Learning Curves of Ensemble Classifiers}
\author{Ada Lovelace}
\academicyear{2025/26}
\studentid{123456}
\keywords{ensemble methods, random forests, learning curves}
\supervisor{Prof.~Super Visor}
\cosupervisor{Prof.~Cosu Pervisor}
\cosupervisor{Dr.~Other Supervisor}
\end{Verbatim}
\cs{academicyear} only takes the years: ``Anno Accademico'' or ``Academic
Year'' is added according to the language. \cs{studentid} is optional; when
present, the student ID is printed below the author's name. \cs{keywords}
is optional too: the keywords are not printed, but stored in the metadata of
the PDF file, together with title, author and degree course, so that search
engines and repositories can index the thesis.
\cs{cosupervisor} can be repeated, or omitted: the label is singular or
plural depending on the number of co-supervisors, and disappears if there
are none. Write \verb|Prof.~| and \verb|Dr.~| with a tie, so that no
sentence-ending space is added after the abbreviation and title and name
stay on the same line.

For the Informatica programmes the cover shows the official logo, which
already contains faculty and course. For the other programmes the faculty
and course are typeset under the university seal; interfaculty programmes
(DSE and DSEH) list both faculties, at the same size. These lines can be
changed with
\begin{Verbatim}
\faculty{Facoltà di Scienze e Tecnologie}
\course{Master's Degree in Data Science\\for Economics and Health}
\end{Verbatim}
The first \cs{faculty} replaces the default faculties, further calls add
lines. A course name that is too long is wrapped automatically; use
\verb|\\| to choose where it breaks.

\subsection{License} \label{sec:license}

A thesis can be released under a license, which tells readers what they are
allowed to do with it. The license is printed at the bottom of the page
following the cover. For Creative Commons licenses, give their type:
\begin{Verbatim}
\license{by-sa}
\end{Verbatim}
The available types are listed in Table~\ref{tab:licenses}; an optional
argument selects a version other than 4.0, as in \verb|\license[3.0]{by}|.
The page shows the official button of the license, a copyright line with
the current year and the author's name (except for CC0, which waives
copyright), and the name and address of the license, in the language of
the document. Any other license or notice can be given as free text:
\begin{Verbatim}
\customlicense{\textcopyright\ 2026 Ada Lovelace. All rights reserved.}
\end{Verbatim}
Without \cs{license} or \cs{customlicense}, the page is omitted.

\begin{table}[htbp]
  \centering
  \caption{Creative Commons licenses available through \cs{license}.}
  \label{tab:licenses}
  \begin{tabular}{ll}
    \toprule
    Type & License \\
    \midrule
    \texttt{by}       & Attribution \\
    \texttt{by-sa}    & Attribution-ShareAlike \\
    \texttt{by-nd}    & Attribution-NoDerivatives \\
    \texttt{by-nc}    & Attribution-NonCommercial \\
    \texttt{by-nc-sa} & Attribution-NonCommercial-ShareAlike \\
    \texttt{by-nc-nd} & Attribution-NonCommercial-NoDerivatives \\
    \texttt{zero}     & CC0 (public domain dedication) \\
    \bottomrule
  \end{tabular}
\end{table}

\section{Structure of the document} \label{sec:structure}

The body of \texttt{main.tex} follows the usual structure of a book:
\begin{Verbatim}
\begin{document}
\frontmatter
\maketitle
\input{frontmatter/dedication}
\input{frontmatter/epigraph}
\input{frontmatter/acknowledgements}
\tableofcontents
\listoffigures
\listoftables
\listofalgorithms
\input{frontmatter/symbols}

\mainmatter
\include{chapters/introduction}
\include{chapters/c1}
...
\include{chapters/conclusions}

\appendix
\include{chapters/appendix}

\printthesisbibliography
\end{document}
\end{Verbatim}
Pages in the front matter are numbered with roman numerals and have no
running header; \cs{mainmatter} restarts the numbering from 1 and switches
on headers and one-and-a-half line spacing; \cs{appendix} starts the
appendices (Appendix~\ref{app:appendices}). The lists of figures, tables
and algorithms, produced by \cs{listoffigures}, \cs{listoftables} and
\cs{listofalgorithms}, collect their captions with the page numbers, and
are themselves listed in the table of contents. They are optional: remove
the corresponding line if the thesis has few figures, tables or
algorithms. A caption that
is too long for these lists can be given a shorter version as optional
argument, which is used in the lists in place of the full caption:
\begin{Verbatim}
\caption[Learning curves]{Learning curves: F1 score on the test set ...}
\end{Verbatim}

\subsection{Front matter}

The dedication and the epigraph have their own commands, which typeset them
on a separate page:
\begin{Verbatim}
\dedication{To my family.}
\epigraphpage{The best way to predict the future\\
          is to invent it.}{--- Alan Kay}
\end{Verbatim}
Both are optional: remove the corresponding \cs{input} line in
\texttt{main.tex} if you do not want them. The acknowledgements are an
unnumbered chapter, described next.

\subsection{List of symbols} \label{sec:symbols}

A thesis with much notation benefits from a list of symbols, placed after
the other lists. It is written by hand in \nolinkurl{frontmatter/symbols.tex},
with one \cs{item} per symbol inside the \texttt{listofsymbols}
environment, which starts an unnumbered chapter titled ``List of Symbols''
(``Elenco dei simboli'' in Italian) and aligns the meanings in a column:
\begin{Verbatim}
\begin{listofsymbols}
  \item[$\R$]                 the set of real numbers
  \item[$\bx_i$]              the $i$-th point of a dataset
  \item[$\lVert \bx \rVert$]  Euclidean norm of $\bx$
\end{listofsymbols}
\end{Verbatim}
Symbols are listed in the order in which they are written, so that related
notation can be kept together. The column of symbols is 2.5\,cm wide; a
different width can be given as optional argument, as in
\verb|\begin{listofsymbols}[4cm]|. If the list is not needed, remove from
\texttt{main.tex} the line that inputs \nolinkurl{frontmatter/symbols.tex}.

\subsection{Unnumbered chapters and sections}

The starred commands \verb|\chapter*|, \verb|\section*| and
\verb|\subsection*| produce headings without number. Unlike standard
\LaTeX{}, they are nevertheless listed in the table of contents, and
\verb|\chapter*| also sets the running header. Unnumbered chapters are the
right choice for the introduction, the acknowledgements and the
conclusions:
\begin{Verbatim}
\chapter*{Acknowledgements}
\section*{Requirements}
\end{Verbatim}
As for numbered headings, an optional argument gives a shorter title for
the table of contents and the header:
\verb|\chapter*[Notation]{Notation and abbreviations}|. An unnumbered
chapter, such as the introduction, should not contain sections: use
\cs{paragraph} headings instead, which start a paragraph and are not listed
in the table of contents, as in the Introduction of this document.

\subsection{Chapters}

Each chapter is a separate file in \texttt{chapters/}, included with
\cs{include}. To add a chapter, create the file and add the corresponding
line to \texttt{main.tex}. To compile only the chapter you are working on,
add \cs{includeonly}\verb|{chapters/c2}| to the preamble: page numbers and
references to the other chapters are preserved.

\subsection{Two-sided printing} \label{sec:twoside}

The template is set up for one-sided printing, which is the most common
choice for a thesis. For a thesis printed on both sides of the sheet,
replace \texttt{oneside} with \texttt{twoside} in the first line of
\texttt{main.tex}:
\begin{Verbatim}
\documentclass[12pt, a4paper, twoside]{book}
\end{Verbatim}
The style then adapts the layout: running headers and page numbers move to
the outer side of the page, on the left of even pages and on the right of
odd ones; chapters, as well as dedication, epigraph and lists, start on a
right-hand page, and the pages left blank for this reason have neither
header nor page number; the license is printed on the back of the cover,
and the ``Made in SLIM'' page on the back cover. If the sheets are bound,
some space can be reserved for the binding on the inner side of the pages,
adding to the preamble, for instance, \verb|\geometry{bindingoffset=1cm}|.

...
\chapter{Before submitting} \label{ch:conclusions}

A few checks are worth doing before the final submission.

\begin{itemize}
  \item The cover shows the right programme, the full names of supervisor
    and co-supervisors, the academic year and, if required, the student ID.
  \item The compilation completes without errors, and the log contains no
    \emph{undefined references} or \emph{undefined citations}: run the full
    sequence of Chapter~\ref{ch1} (or \texttt{latexmk}) once more after the
    last change.
  \item The log reports no \emph{overfull boxes}, i.e., no lines exceeding
    the margin; long words or formulas causing them can be fixed by
    rephrasing, or by suggesting hyphenation points with \verb|\-|.
  \item Every figure and table is referred to in the text, and its caption
    can be understood without reading the surrounding text.
  \item Bibliographic entries are complete (authors, title, venue, year) and
    have a DOI whenever one exists.
  \item Placeholder material, such as the examples of this document, the
    dedication or the epigraph if not wanted, has been removed.
\end{itemize}

\section*{Customising the style}

The style is meant to be used as it is, but a few aspects are easy to
change in \texttt{slimthesis.sty}, where they are grouped and commented:
the spacing of the bibliography (\cs{bibitemsep} and the line spacing set by
\cs{AtBeginBibliography}), the colour of links (\texttt{slimblue}), and the
fonts. Keep such changes to a minimum, and note them down: when a new
version of the style is released, they will have to be applied again.

The last page of the thesis acknowledges the SLIM laboratory; it can be
omitted with the \texttt{madeinslim=false} option described in
Section~\ref{sec:programme}.


\appendix
\chapter{Appendices} \label{app:appendices}

Material that supports the thesis but would interrupt its flow, such as
long proofs, additional experimental results, the full description of a
dataset or a user manual, belongs in one or more appendices, placed after
the conclusions and before the bibliography.

\section{Writing an appendix}

The command \cs{appendix} in \texttt{main.tex} marks the beginning of the
appendices; each of them is then an ordinary chapter, in its own file:
\begin{Verbatim}
\include{chapters/conclusions}

\appendix
\include{chapters/appendix}

\printthesisbibliography
\end{Verbatim}
After \cs{appendix}, chapters are labelled ``Appendix'' (``Appendice'' in
Italian) instead of ``Chapter'', both in their heading and in the running
header, and are numbered with capital letters: this appendix is
Appendix~\ref{app:appendices}, its first section is
Section~\ref{app:appendices}.1, and its figures, tables, theorems and
algorithms are numbered accordingly, as for Table~\ref{tab:search-space}.
Cross-references work as in the rest of the thesis.

\section{An example of supporting material}

Table~\ref{tab:search-space} lists the values of the hyperparameters
explored by the grid search of Table~\ref{tab:grid}: a typical detail that
is needed to reproduce the experiments, but not to follow the main
discussion.

\begin{table}[htbp]
  \centering
  \caption[Search space of the grid search]{Values of the hyperparameters
    explored by the grid search of Table~\ref{tab:grid}.}
  \label{tab:search-space}
  \begin{tabular}{ll}
    \toprule
    Hyperparameter & Values \\
    \midrule
    Maximum depth of the trees      & 4, 8, 16, 32, \texttt{None} (unbounded) \\
    Number of trees                 & 50, 100, 200, 500 \\
    Features considered at a split  & \texttt{sqrt}, \texttt{log2}, 0.5 \\
    \bottomrule
  \end{tabular}
\end{table}

An appendix should be referred to from the main text, at the point where
its content is relevant; material that is never referred to probably does
not belong in the thesis.


\printthesisbibliography
\end{document}
\end{Verbatim}
Pages in the front matter are numbered with roman numerals and have no
running header; \cs{mainmatter} restarts the numbering from 1 and switches
on headers and one-and-a-half line spacing; \cs{appendix} starts the
appendices (Appendix~\ref{app:appendices}). The lists of figures, tables
and algorithms, produced by \cs{listoffigures}, \cs{listoftables} and
\cs{listofalgorithms}, collect their captions with the page numbers, and
are themselves listed in the table of contents. They are optional: remove
the corresponding line if the thesis has few figures, tables or
algorithms. A caption that
is too long for these lists can be given a shorter version as optional
argument, which is used in the lists in place of the full caption:
\begin{Verbatim}
\caption[Learning curves]{Learning curves: F1 score on the test set ...}
\end{Verbatim}

\subsection{Front matter}

The dedication and the epigraph have their own commands, which typeset them
on a separate page:
\begin{Verbatim}
\dedication{To my family.}
\epigraphpage{The best way to predict the future\\
          is to invent it.}{--- Alan Kay}
\end{Verbatim}
Both are optional: remove the corresponding \cs{input} line in
\texttt{main.tex} if you do not want them. The acknowledgements are an
unnumbered chapter, described next.

\subsection{List of symbols} \label{sec:symbols}

A thesis with much notation benefits from a list of symbols, placed after
the other lists. It is written by hand in \nolinkurl{frontmatter/symbols.tex},
with one \cs{item} per symbol inside the \texttt{listofsymbols}
environment, which starts an unnumbered chapter titled ``List of Symbols''
(``Elenco dei simboli'' in Italian) and aligns the meanings in a column:
\begin{Verbatim}
\begin{listofsymbols}
  \item[$\R$]                 the set of real numbers
  \item[$\bx_i$]              the $i$-th point of a dataset
  \item[$\lVert \bx \rVert$]  Euclidean norm of $\bx$
\end{listofsymbols}
\end{Verbatim}
Symbols are listed in the order in which they are written, so that related
notation can be kept together. The column of symbols is 2.5\,cm wide; a
different width can be given as optional argument, as in
\verb|\begin{listofsymbols}[4cm]|. If the list is not needed, remove from
\texttt{main.tex} the line that inputs \nolinkurl{frontmatter/symbols.tex}.

\subsection{Unnumbered chapters and sections}

The starred commands \verb|\chapter*|, \verb|\section*| and
\verb|\subsection*| produce headings without number. Unlike standard
\LaTeX{}, they are nevertheless listed in the table of contents, and
\verb|\chapter*| also sets the running header. Unnumbered chapters are the
right choice for the introduction, the acknowledgements and the
conclusions:
\begin{Verbatim}
\chapter*{Acknowledgements}
\section*{Requirements}
\end{Verbatim}
As for numbered headings, an optional argument gives a shorter title for
the table of contents and the header:
\verb|\chapter*[Notation]{Notation and abbreviations}|. An unnumbered
chapter, such as the introduction, should not contain sections: use
\cs{paragraph} headings instead, which start a paragraph and are not listed
in the table of contents, as in the Introduction of this document.

\subsection{Chapters}

Each chapter is a separate file in \texttt{chapters/}, included with
\cs{include}. To add a chapter, create the file and add the corresponding
line to \texttt{main.tex}. To compile only the chapter you are working on,
add \cs{includeonly}\verb|{chapters/c2}| to the preamble: page numbers and
references to the other chapters are preserved.

\subsection{Two-sided printing} \label{sec:twoside}

The template is set up for one-sided printing, which is the most common
choice for a thesis. For a thesis printed on both sides of the sheet,
replace \texttt{oneside} with \texttt{twoside} in the first line of
\texttt{main.tex}:
\begin{Verbatim}
\documentclass[12pt, a4paper, twoside]{book}
\end{Verbatim}
The style then adapts the layout: running headers and page numbers move to
the outer side of the page, on the left of even pages and on the right of
odd ones; chapters, as well as dedication, epigraph and lists, start on a
right-hand page, and the pages left blank for this reason have neither
header nor page number; the license is printed on the back of the cover,
and the ``Made in SLIM'' page on the back cover. If the sheets are bound,
some space can be reserved for the binding on the inner side of the pages,
adding to the preamble, for instance, \verb|\geometry{bindingoffset=1cm}|.

...
\chapter{Before submitting} \label{ch:conclusions}

A few checks are worth doing before the final submission.

\begin{itemize}
  \item The cover shows the right programme, the full names of supervisor
    and co-supervisors, the academic year and, if required, the student ID.
  \item The compilation completes without errors, and the log contains no
    \emph{undefined references} or \emph{undefined citations}: run the full
    sequence of Chapter~\ref{ch1} (or \texttt{latexmk}) once more after the
    last change.
  \item The log reports no \emph{overfull boxes}, i.e., no lines exceeding
    the margin; long words or formulas causing them can be fixed by
    rephrasing, or by suggesting hyphenation points with \verb|\-|.
  \item Every figure and table is referred to in the text, and its caption
    can be understood without reading the surrounding text.
  \item Bibliographic entries are complete (authors, title, venue, year) and
    have a DOI whenever one exists.
  \item Placeholder material, such as the examples of this document, the
    dedication or the epigraph if not wanted, has been removed.
\end{itemize}

\section*{Customising the style}

The style is meant to be used as it is, but a few aspects are easy to
change in \texttt{slimthesis.sty}, where they are grouped and commented:
the spacing of the bibliography (\cs{bibitemsep} and the line spacing set by
\cs{AtBeginBibliography}), the colour of links (\texttt{slimblue}), and the
fonts. Keep such changes to a minimum, and note them down: when a new
version of the style is released, they will have to be applied again.

The last page of the thesis acknowledges the SLIM laboratory; it can be
omitted with the \texttt{madeinslim=false} option described in
Section~\ref{sec:programme}.


\appendix
\chapter{Appendices} \label{app:appendices}

Material that supports the thesis but would interrupt its flow, such as
long proofs, additional experimental results, the full description of a
dataset or a user manual, belongs in one or more appendices, placed after
the conclusions and before the bibliography.

\section{Writing an appendix}

The command \cs{appendix} in \texttt{main.tex} marks the beginning of the
appendices; each of them is then an ordinary chapter, in its own file:
\begin{Verbatim}
\chapter{Before submitting} \label{ch:conclusions}

A few checks are worth doing before the final submission.

\begin{itemize}
  \item The cover shows the right programme, the full names of supervisor
    and co-supervisors, the academic year and, if required, the student ID.
  \item The compilation completes without errors, and the log contains no
    \emph{undefined references} or \emph{undefined citations}: run the full
    sequence of Chapter~\ref{ch1} (or \texttt{latexmk}) once more after the
    last change.
  \item The log reports no \emph{overfull boxes}, i.e., no lines exceeding
    the margin; long words or formulas causing them can be fixed by
    rephrasing, or by suggesting hyphenation points with \verb|\-|.
  \item Every figure and table is referred to in the text, and its caption
    can be understood without reading the surrounding text.
  \item Bibliographic entries are complete (authors, title, venue, year) and
    have a DOI whenever one exists.
  \item Placeholder material, such as the examples of this document, the
    dedication or the epigraph if not wanted, has been removed.
\end{itemize}

\section*{Customising the style}

The style is meant to be used as it is, but a few aspects are easy to
change in \texttt{slimthesis.sty}, where they are grouped and commented:
the spacing of the bibliography (\cs{bibitemsep} and the line spacing set by
\cs{AtBeginBibliography}), the colour of links (\texttt{slimblue}), and the
fonts. Keep such changes to a minimum, and note them down: when a new
version of the style is released, they will have to be applied again.

The last page of the thesis acknowledges the SLIM laboratory; it can be
omitted with the \texttt{madeinslim=false} option described in
Section~\ref{sec:programme}.


\appendix
\chapter{Appendices} \label{app:appendices}

Material that supports the thesis but would interrupt its flow, such as
long proofs, additional experimental results, the full description of a
dataset or a user manual, belongs in one or more appendices, placed after
the conclusions and before the bibliography.

\section{Writing an appendix}

The command \cs{appendix} in \texttt{main.tex} marks the beginning of the
appendices; each of them is then an ordinary chapter, in its own file:
\begin{Verbatim}
\include{chapters/conclusions}

\appendix
\include{chapters/appendix}

\printthesisbibliography
\end{Verbatim}
After \cs{appendix}, chapters are labelled ``Appendix'' (``Appendice'' in
Italian) instead of ``Chapter'', both in their heading and in the running
header, and are numbered with capital letters: this appendix is
Appendix~\ref{app:appendices}, its first section is
Section~\ref{app:appendices}.1, and its figures, tables, theorems and
algorithms are numbered accordingly, as for Table~\ref{tab:search-space}.
Cross-references work as in the rest of the thesis.

\section{An example of supporting material}

Table~\ref{tab:search-space} lists the values of the hyperparameters
explored by the grid search of Table~\ref{tab:grid}: a typical detail that
is needed to reproduce the experiments, but not to follow the main
discussion.

\begin{table}[htbp]
  \centering
  \caption[Search space of the grid search]{Values of the hyperparameters
    explored by the grid search of Table~\ref{tab:grid}.}
  \label{tab:search-space}
  \begin{tabular}{ll}
    \toprule
    Hyperparameter & Values \\
    \midrule
    Maximum depth of the trees      & 4, 8, 16, 32, \texttt{None} (unbounded) \\
    Number of trees                 & 50, 100, 200, 500 \\
    Features considered at a split  & \texttt{sqrt}, \texttt{log2}, 0.5 \\
    \bottomrule
  \end{tabular}
\end{table}

An appendix should be referred to from the main text, at the point where
its content is relevant; material that is never referred to probably does
not belong in the thesis.


\printthesisbibliography
\end{Verbatim}
After \cs{appendix}, chapters are labelled ``Appendix'' (``Appendice'' in
Italian) instead of ``Chapter'', both in their heading and in the running
header, and are numbered with capital letters: this appendix is
Appendix~\ref{app:appendices}, its first section is
Section~\ref{app:appendices}.1, and its figures, tables, theorems and
algorithms are numbered accordingly, as for Table~\ref{tab:search-space}.
Cross-references work as in the rest of the thesis.

\section{An example of supporting material}

Table~\ref{tab:search-space} lists the values of the hyperparameters
explored by the grid search of Table~\ref{tab:grid}: a typical detail that
is needed to reproduce the experiments, but not to follow the main
discussion.

\begin{table}[htbp]
  \centering
  \caption[Search space of the grid search]{Values of the hyperparameters
    explored by the grid search of Table~\ref{tab:grid}.}
  \label{tab:search-space}
  \begin{tabular}{ll}
    \toprule
    Hyperparameter & Values \\
    \midrule
    Maximum depth of the trees      & 4, 8, 16, 32, \texttt{None} (unbounded) \\
    Number of trees                 & 50, 100, 200, 500 \\
    Features considered at a split  & \texttt{sqrt}, \texttt{log2}, 0.5 \\
    \bottomrule
  \end{tabular}
\end{table}

An appendix should be referred to from the main text, at the point where
its content is relevant; material that is never referred to probably does
not belong in the thesis.


\printthesisbibliography
\end{document}
\end{Verbatim}
Pages in the front matter are numbered with roman numerals and have no
running header; \cs{mainmatter} restarts the numbering from 1 and switches
on headers and one-and-a-half line spacing; \cs{appendix} starts the
appendices (Appendix~\ref{app:appendices}). The lists of figures, tables
and algorithms, produced by \cs{listoffigures}, \cs{listoftables} and
\cs{listofalgorithms}, collect their captions with the page numbers, and
are themselves listed in the table of contents. They are optional: remove
the corresponding line if the thesis has few figures, tables or
algorithms. A caption that
is too long for these lists can be given a shorter version as optional
argument, which is used in the lists in place of the full caption:
\begin{Verbatim}
\caption[Learning curves]{Learning curves: F1 score on the test set ...}
\end{Verbatim}

\subsection{Front matter}

The dedication and the epigraph have their own commands, which typeset them
on a separate page:
\begin{Verbatim}
\dedication{To my family.}
\epigraphpage{The best way to predict the future\\
          is to invent it.}{--- Alan Kay}
\end{Verbatim}
Both are optional: remove the corresponding \cs{input} line in
\texttt{main.tex} if you do not want them. The acknowledgements are an
unnumbered chapter, described next.

\subsection{List of symbols} \label{sec:symbols}

A thesis with much notation benefits from a list of symbols, placed after
the other lists. It is written by hand in \nolinkurl{frontmatter/symbols.tex},
with one \cs{item} per symbol inside the \texttt{listofsymbols}
environment, which starts an unnumbered chapter titled ``List of Symbols''
(``Elenco dei simboli'' in Italian) and aligns the meanings in a column:
\begin{Verbatim}
\begin{listofsymbols}
  \item[$\R$]                 the set of real numbers
  \item[$\bx_i$]              the $i$-th point of a dataset
  \item[$\lVert \bx \rVert$]  Euclidean norm of $\bx$
\end{listofsymbols}
\end{Verbatim}
Symbols are listed in the order in which they are written, so that related
notation can be kept together. The column of symbols is 2.5\,cm wide; a
different width can be given as optional argument, as in
\verb|\begin{listofsymbols}[4cm]|. If the list is not needed, remove from
\texttt{main.tex} the line that inputs \nolinkurl{frontmatter/symbols.tex}.

\subsection{Unnumbered chapters and sections}

The starred commands \verb|\chapter*|, \verb|\section*| and
\verb|\subsection*| produce headings without number. Unlike standard
\LaTeX{}, they are nevertheless listed in the table of contents, and
\verb|\chapter*| also sets the running header. Unnumbered chapters are the
right choice for the introduction, the acknowledgements and the
conclusions:
\begin{Verbatim}
\chapter*{Acknowledgements}
\section*{Requirements}
\end{Verbatim}
As for numbered headings, an optional argument gives a shorter title for
the table of contents and the header:
\verb|\chapter*[Notation]{Notation and abbreviations}|. An unnumbered
chapter, such as the introduction, should not contain sections: use
\cs{paragraph} headings instead, which start a paragraph and are not listed
in the table of contents, as in the Introduction of this document.

\subsection{Chapters}

Each chapter is a separate file in \texttt{chapters/}, included with
\cs{include}. To add a chapter, create the file and add the corresponding
line to \texttt{main.tex}. To compile only the chapter you are working on,
add \cs{includeonly}\verb|{chapters/c2}| to the preamble: page numbers and
references to the other chapters are preserved.

\subsection{Two-sided printing} \label{sec:twoside}

The template is set up for one-sided printing, which is the most common
choice for a thesis. For a thesis printed on both sides of the sheet,
replace \texttt{oneside} with \texttt{twoside} in the first line of
\texttt{main.tex}:
\begin{Verbatim}
\documentclass[12pt, a4paper, twoside]{book}
\end{Verbatim}
The style then adapts the layout: running headers and page numbers move to
the outer side of the page, on the left of even pages and on the right of
odd ones; chapters, as well as dedication, epigraph and lists, start on a
right-hand page, and the pages left blank for this reason have neither
header nor page number; the license is printed on the back of the cover,
and the ``Made in SLIM'' page on the back cover. If the sheets are bound,
some space can be reserved for the binding on the inner side of the pages,
adding to the preamble, for instance, \verb|\geometry{bindingoffset=1cm}|.

...
\chapter{Before submitting} \label{ch:conclusions}

A few checks are worth doing before the final submission.

\begin{itemize}
  \item The cover shows the right programme, the full names of supervisor
    and co-supervisors, the academic year and, if required, the student ID.
  \item The compilation completes without errors, and the log contains no
    \emph{undefined references} or \emph{undefined citations}: run the full
    sequence of Chapter~\ref{ch1} (or \texttt{latexmk}) once more after the
    last change.
  \item The log reports no \emph{overfull boxes}, i.e., no lines exceeding
    the margin; long words or formulas causing them can be fixed by
    rephrasing, or by suggesting hyphenation points with \verb|\-|.
  \item Every figure and table is referred to in the text, and its caption
    can be understood without reading the surrounding text.
  \item Bibliographic entries are complete (authors, title, venue, year) and
    have a DOI whenever one exists.
  \item Placeholder material, such as the examples of this document, the
    dedication or the epigraph if not wanted, has been removed.
\end{itemize}

\section*{Customising the style}

The style is meant to be used as it is, but a few aspects are easy to
change in \texttt{slimthesis.sty}, where they are grouped and commented:
the spacing of the bibliography (\cs{bibitemsep} and the line spacing set by
\cs{AtBeginBibliography}), the colour of links (\texttt{slimblue}), and the
fonts. Keep such changes to a minimum, and note them down: when a new
version of the style is released, they will have to be applied again.

The last page of the thesis acknowledges the SLIM laboratory; it can be
omitted with the \texttt{madeinslim=false} option described in
Section~\ref{sec:programme}.


\appendix
\chapter{Appendices} \label{app:appendices}

Material that supports the thesis but would interrupt its flow, such as
long proofs, additional experimental results, the full description of a
dataset or a user manual, belongs in one or more appendices, placed after
the conclusions and before the bibliography.

\section{Writing an appendix}

The command \cs{appendix} in \texttt{main.tex} marks the beginning of the
appendices; each of them is then an ordinary chapter, in its own file:
\begin{Verbatim}
\chapter{Before submitting} \label{ch:conclusions}

A few checks are worth doing before the final submission.

\begin{itemize}
  \item The cover shows the right programme, the full names of supervisor
    and co-supervisors, the academic year and, if required, the student ID.
  \item The compilation completes without errors, and the log contains no
    \emph{undefined references} or \emph{undefined citations}: run the full
    sequence of Chapter~\ref{ch1} (or \texttt{latexmk}) once more after the
    last change.
  \item The log reports no \emph{overfull boxes}, i.e., no lines exceeding
    the margin; long words or formulas causing them can be fixed by
    rephrasing, or by suggesting hyphenation points with \verb|\-|.
  \item Every figure and table is referred to in the text, and its caption
    can be understood without reading the surrounding text.
  \item Bibliographic entries are complete (authors, title, venue, year) and
    have a DOI whenever one exists.
  \item Placeholder material, such as the examples of this document, the
    dedication or the epigraph if not wanted, has been removed.
\end{itemize}

\section*{Customising the style}

The style is meant to be used as it is, but a few aspects are easy to
change in \texttt{slimthesis.sty}, where they are grouped and commented:
the spacing of the bibliography (\cs{bibitemsep} and the line spacing set by
\cs{AtBeginBibliography}), the colour of links (\texttt{slimblue}), and the
fonts. Keep such changes to a minimum, and note them down: when a new
version of the style is released, they will have to be applied again.

The last page of the thesis acknowledges the SLIM laboratory; it can be
omitted with the \texttt{madeinslim=false} option described in
Section~\ref{sec:programme}.


\appendix
\chapter{Appendices} \label{app:appendices}

Material that supports the thesis but would interrupt its flow, such as
long proofs, additional experimental results, the full description of a
dataset or a user manual, belongs in one or more appendices, placed after
the conclusions and before the bibliography.

\section{Writing an appendix}

The command \cs{appendix} in \texttt{main.tex} marks the beginning of the
appendices; each of them is then an ordinary chapter, in its own file:
\begin{Verbatim}
\chapter{Before submitting} \label{ch:conclusions}

A few checks are worth doing before the final submission.

\begin{itemize}
  \item The cover shows the right programme, the full names of supervisor
    and co-supervisors, the academic year and, if required, the student ID.
  \item The compilation completes without errors, and the log contains no
    \emph{undefined references} or \emph{undefined citations}: run the full
    sequence of Chapter~\ref{ch1} (or \texttt{latexmk}) once more after the
    last change.
  \item The log reports no \emph{overfull boxes}, i.e., no lines exceeding
    the margin; long words or formulas causing them can be fixed by
    rephrasing, or by suggesting hyphenation points with \verb|\-|.
  \item Every figure and table is referred to in the text, and its caption
    can be understood without reading the surrounding text.
  \item Bibliographic entries are complete (authors, title, venue, year) and
    have a DOI whenever one exists.
  \item Placeholder material, such as the examples of this document, the
    dedication or the epigraph if not wanted, has been removed.
\end{itemize}

\section*{Customising the style}

The style is meant to be used as it is, but a few aspects are easy to
change in \texttt{slimthesis.sty}, where they are grouped and commented:
the spacing of the bibliography (\cs{bibitemsep} and the line spacing set by
\cs{AtBeginBibliography}), the colour of links (\texttt{slimblue}), and the
fonts. Keep such changes to a minimum, and note them down: when a new
version of the style is released, they will have to be applied again.

The last page of the thesis acknowledges the SLIM laboratory; it can be
omitted with the \texttt{madeinslim=false} option described in
Section~\ref{sec:programme}.


\appendix
\chapter{Appendices} \label{app:appendices}

Material that supports the thesis but would interrupt its flow, such as
long proofs, additional experimental results, the full description of a
dataset or a user manual, belongs in one or more appendices, placed after
the conclusions and before the bibliography.

\section{Writing an appendix}

The command \cs{appendix} in \texttt{main.tex} marks the beginning of the
appendices; each of them is then an ordinary chapter, in its own file:
\begin{Verbatim}
\include{chapters/conclusions}

\appendix
\include{chapters/appendix}

\printthesisbibliography
\end{Verbatim}
After \cs{appendix}, chapters are labelled ``Appendix'' (``Appendice'' in
Italian) instead of ``Chapter'', both in their heading and in the running
header, and are numbered with capital letters: this appendix is
Appendix~\ref{app:appendices}, its first section is
Section~\ref{app:appendices}.1, and its figures, tables, theorems and
algorithms are numbered accordingly, as for Table~\ref{tab:search-space}.
Cross-references work as in the rest of the thesis.

\section{An example of supporting material}

Table~\ref{tab:search-space} lists the values of the hyperparameters
explored by the grid search of Table~\ref{tab:grid}: a typical detail that
is needed to reproduce the experiments, but not to follow the main
discussion.

\begin{table}[htbp]
  \centering
  \caption[Search space of the grid search]{Values of the hyperparameters
    explored by the grid search of Table~\ref{tab:grid}.}
  \label{tab:search-space}
  \begin{tabular}{ll}
    \toprule
    Hyperparameter & Values \\
    \midrule
    Maximum depth of the trees      & 4, 8, 16, 32, \texttt{None} (unbounded) \\
    Number of trees                 & 50, 100, 200, 500 \\
    Features considered at a split  & \texttt{sqrt}, \texttt{log2}, 0.5 \\
    \bottomrule
  \end{tabular}
\end{table}

An appendix should be referred to from the main text, at the point where
its content is relevant; material that is never referred to probably does
not belong in the thesis.


\printthesisbibliography
\end{Verbatim}
After \cs{appendix}, chapters are labelled ``Appendix'' (``Appendice'' in
Italian) instead of ``Chapter'', both in their heading and in the running
header, and are numbered with capital letters: this appendix is
Appendix~\ref{app:appendices}, its first section is
Section~\ref{app:appendices}.1, and its figures, tables, theorems and
algorithms are numbered accordingly, as for Table~\ref{tab:search-space}.
Cross-references work as in the rest of the thesis.

\section{An example of supporting material}

Table~\ref{tab:search-space} lists the values of the hyperparameters
explored by the grid search of Table~\ref{tab:grid}: a typical detail that
is needed to reproduce the experiments, but not to follow the main
discussion.

\begin{table}[htbp]
  \centering
  \caption[Search space of the grid search]{Values of the hyperparameters
    explored by the grid search of Table~\ref{tab:grid}.}
  \label{tab:search-space}
  \begin{tabular}{ll}
    \toprule
    Hyperparameter & Values \\
    \midrule
    Maximum depth of the trees      & 4, 8, 16, 32, \texttt{None} (unbounded) \\
    Number of trees                 & 50, 100, 200, 500 \\
    Features considered at a split  & \texttt{sqrt}, \texttt{log2}, 0.5 \\
    \bottomrule
  \end{tabular}
\end{table}

An appendix should be referred to from the main text, at the point where
its content is relevant; material that is never referred to probably does
not belong in the thesis.


\printthesisbibliography
\end{Verbatim}
After \cs{appendix}, chapters are labelled ``Appendix'' (``Appendice'' in
Italian) instead of ``Chapter'', both in their heading and in the running
header, and are numbered with capital letters: this appendix is
Appendix~\ref{app:appendices}, its first section is
Section~\ref{app:appendices}.1, and its figures, tables, theorems and
algorithms are numbered accordingly, as for Table~\ref{tab:search-space}.
Cross-references work as in the rest of the thesis.

\section{An example of supporting material}

Table~\ref{tab:search-space} lists the values of the hyperparameters
explored by the grid search of Table~\ref{tab:grid}: a typical detail that
is needed to reproduce the experiments, but not to follow the main
discussion.

\begin{table}[htbp]
  \centering
  \caption[Search space of the grid search]{Values of the hyperparameters
    explored by the grid search of Table~\ref{tab:grid}.}
  \label{tab:search-space}
  \begin{tabular}{ll}
    \toprule
    Hyperparameter & Values \\
    \midrule
    Maximum depth of the trees      & 4, 8, 16, 32, \texttt{None} (unbounded) \\
    Number of trees                 & 50, 100, 200, 500 \\
    Features considered at a split  & \texttt{sqrt}, \texttt{log2}, 0.5 \\
    \bottomrule
  \end{tabular}
\end{table}

An appendix should be referred to from the main text, at the point where
its content is relevant; material that is never referred to probably does
not belong in the thesis.


\printthesisbibliography
\end{document}
\end{Verbatim}
Pages in the front matter are numbered with roman numerals and have no
running header; \cs{mainmatter} restarts the numbering from 1 and switches
on headers and one-and-a-half line spacing; \cs{appendix} starts the
appendices (Appendix~\ref{app:appendices}). The lists of figures, tables
and algorithms, produced by \cs{listoffigures}, \cs{listoftables} and
\cs{listofalgorithms}, collect their captions with the page numbers, and
are themselves listed in the table of contents. They are optional: remove
the corresponding line if the thesis has few figures, tables or
algorithms. A caption that
is too long for these lists can be given a shorter version as optional
argument, which is used in the lists in place of the full caption:
\begin{Verbatim}
\caption[Learning curves]{Learning curves: F1 score on the test set ...}
\end{Verbatim}

\subsection{Front matter}

The dedication and the epigraph have their own commands, which typeset them
on a separate page:
\begin{Verbatim}
\dedication{To my family.}
\epigraphpage{The best way to predict the future\\
          is to invent it.}{--- Alan Kay}
\end{Verbatim}
Both are optional: remove the corresponding \cs{input} line in
\texttt{main.tex} if you do not want them. The acknowledgements are an
unnumbered chapter, described next.

\subsection{List of symbols} \label{sec:symbols}

A thesis with much notation benefits from a list of symbols, placed after
the other lists. It is written by hand in \nolinkurl{frontmatter/symbols.tex},
with one \cs{item} per symbol inside the \texttt{listofsymbols}
environment, which starts an unnumbered chapter titled ``List of Symbols''
(``Elenco dei simboli'' in Italian) and aligns the meanings in a column:
\begin{Verbatim}
\begin{listofsymbols}
  \item[$\R$]                 the set of real numbers
  \item[$\bx_i$]              the $i$-th point of a dataset
  \item[$\lVert \bx \rVert$]  Euclidean norm of $\bx$
\end{listofsymbols}
\end{Verbatim}
Symbols are listed in the order in which they are written, so that related
notation can be kept together. The column of symbols is 2.5\,cm wide; a
different width can be given as optional argument, as in
\verb|\begin{listofsymbols}[4cm]|. If the list is not needed, remove from
\texttt{main.tex} the line that inputs \nolinkurl{frontmatter/symbols.tex}.

\subsection{Unnumbered chapters and sections}

The starred commands \verb|\chapter*|, \verb|\section*| and
\verb|\subsection*| produce headings without number. Unlike standard
\LaTeX{}, they are nevertheless listed in the table of contents, and
\verb|\chapter*| also sets the running header. Unnumbered chapters are the
right choice for the introduction, the acknowledgements and the
conclusions:
\begin{Verbatim}
\chapter*{Acknowledgements}
\section*{Requirements}
\end{Verbatim}
As for numbered headings, an optional argument gives a shorter title for
the table of contents and the header:
\verb|\chapter*[Notation]{Notation and abbreviations}|. An unnumbered
chapter, such as the introduction, should not contain sections: use
\cs{paragraph} headings instead, which start a paragraph and are not listed
in the table of contents, as in the Introduction of this document.

\subsection{Chapters}

Each chapter is a separate file in \texttt{chapters/}, included with
\cs{include}. To add a chapter, create the file and add the corresponding
line to \texttt{main.tex}. To compile only the chapter you are working on,
add \cs{includeonly}\verb|{chapters/c2}| to the preamble: page numbers and
references to the other chapters are preserved.

\subsection{Two-sided printing} \label{sec:twoside}

The template is set up for one-sided printing, which is the most common
choice for a thesis. For a thesis printed on both sides of the sheet,
replace \texttt{oneside} with \texttt{twoside} in the first line of
\texttt{main.tex}:
\begin{Verbatim}
\documentclass[12pt, a4paper, twoside]{book}
\end{Verbatim}
The style then adapts the layout: running headers and page numbers move to
the outer side of the page, on the left of even pages and on the right of
odd ones; chapters, as well as dedication, epigraph and lists, start on a
right-hand page, and the pages left blank for this reason have neither
header nor page number; the license is printed on the back of the cover,
and the ``Made in SLIM'' page on the back cover. If the sheets are bound,
some space can be reserved for the binding on the inner side of the pages,
adding to the preamble, for instance, \verb|\geometry{bindingoffset=1cm}|.
